\documentclass{amsart}

\usepackage{keytheorems}
\usepackage{amsthm}
\usepackage{color}

\newkeytheorem{problem}
\newtheorem*{proposition}{Proposition}
\newcommand{\todo}{\textcolor{red}{TO-DO.} }
\newcommand{\temp}{\textcolor{red}{TEMPORARY.} }

\newcommand{\set}[1]{\left\{#1\right\}}
\newcommand{\powset}{{\cal P}}
\newcommand{\img}{_{*}}
\newcommand{\pimg}{^{*}}
\newcommand{\id}[1]{\mathrm{id}_{#1}}

\newcommand{\R}{\mathbb{R}}
\newcommand{\Z}{\mathbb{Z}}
\newcommand{\C}{\mathbb{C}}
\newcommand{\N}{\mathbb{N}}
\newcommand{\Q}{\mathbb{Q}}

\newcommand{\W}{\mathcal{W}}
\newcommand{\T}{\mathcal{T}}

\numberwithin{equation}{section}

\begin{document}

	\title[Math 251W: Homework \#10]{
		Math 251W: Foundations of Advanced Mathematics \\
		\textmd{Homework \#10: Portfolio Problems \S 5.2 \& 5.3}}
	\author{Carmen Beltran \and Brendan Butler \and Teddy Wachtler}
	\date{\today}

	\maketitle

	\begin{problem}[manual-num=55]
		\temp
		\begin{proposition}
			Let $n, q \in \N$, and let $a, b \in \Z$.
			Suppose $a \equiv b (\bmod{n}),$ and $q|n$.
			Prove that $a \equiv b (\bmod{q})$.
		\end{proposition}

		\begin{proof}
			Let $n$ and $q$ be arbitrary natural numbers.
			Let $a$ and $b$ be arbitrary integers.
			Suppose $a \equiv b (\bmod{n})$ and $q | n$.
			By the definition of modular congruence, we see that there exists a particular integer $c$ so that $a - b = cn$.
			By the definition of divides, we see that there exists a particular integer $d$ so that $n = qd$.
			By substitution, we see that $a - b = cqd$.
			By the closure of the integers, there exists a particular integer $e$ so that $e = cd$.
			By substitution, we see that $a - b = eq$.
			By the definition of modular congruence, we see that $a \equiv b (\bmod{q})$.
			Therefore, we see that for all natural numbers $n$ and $q$, and for all integers $a$ and $b$, if $a \equiv b(\bmod{n})$ and $q|n$, then $a \equiv b (\bmod{q})$.
		\end{proof}
	\end{problem}

	\begin{problem}[manual-num=56]
		Let $A$ be a non-empty set.
		Let $E(A)$ denote the set of all equivalence relations on the set $A$.
		Let $\T_A$ denote the set of all partitions of the set $A$.
		Let $F: E(A) \rightarrow \T_A$ be the function defined by $F(\sim) = A / \sim$ for all equivalence relations $\sim \in E(A)$.
		Let $G : \T_A \rightarrow E(A)$ be the function defined by $G(\W) = R_\W$ for all partitions $\W \in \T_A$.

		\begin{enumerate}
			\item
			Prove the function $F$ is well-defined by showing $F(\sim) \in \T_A$.

			\begin{proof}
				We ask whether $F : E(A) \rightarrow \T_A$ is well-defined.
				Let $\sim$ be an arbitrary element in $E(A)$.
				We ask whether $A / \sim$ forms a partition of $A$.

				Let $P_a$ be an arbitrary set in $A / \sim$.
				By the definition of set quotients, we see that there exists a particular element $a \in A$ so that $P_a = \sim [a]$.
				By the definition of reflexive relations, we see that $a \in \sim[a]$.
				By the substitution, we see that $a \in P_a$.
				By the definition of the empty set, we see that $P_a \neq \emptyset$.
				Since $P_a \in A / \sim$ implies $P_a \neq \emptyset$, we see that $A / \sim$ contains no empty subsets.

				Let $P_b$ and $Q_b$ be arbitrary sets so that $P_b \in A / \sim$ and $Q_b \in A / \sim$.
				Suppose that $P_b \cap Q_b \neq \emptyset$.
				Let $b_0$ be an arbitrary element of $P_b$.
				Let $b_1$ be an arbitrary element of $P_b \cap Q_b$.
				By the definition of set intersection, we see that $b_1 \in P_b$ and $b_1 \in Q_b$.
				By the definition of set quotients, since $P_b \in A / \sim$ and $Q_b \in A / \sim$, we see that there exists particular elements $b_p \in A$ and $b_q \in A$ so that $P_b = \sim [b_p]$ and $Q_b = \sim [b_q]$.
				By substitution, we see that $b_0 \in \sim [b_p]$, $b_1 \in \sim [b_p]$, and $b_1 \in \sim [b_q]$.
				By the definition of equivalence classes, we see that $(b_p, b_0) \in \sim$, $(b_p, b_1) \in \sim$, and $(b_q, b_1) \in \sim$.
				By the definition of symmetric relations, we see that $(b_1, b_p) \in \sim$.
				By the definition of transitive relations, we see that $(b_q, b_p) \in \sim$.
				By the definition of transitive relations, we see that $(b_q, b_0) \in \sim$.
				By the definition of equivalence classes, we see that $b_0 \in \sim [b_q]$.
				By the definition of set quotients, we see that $b_0 \in Q_b$.
				By the definition of subsets, since $b_0 \in P_b$ implies $b_0 \in Q_b$, we see that $P_b \subseteq Q_b$.
				Analogously, we see that $Q_b \subseteq P_b$.
				By the definition of set equality, we see that $P_b = Q_b$.
				By the contrapositive, since $P_b \cap Q_b \neq \emptyset$ implies $P_b = Q_b$, $P_b \neq Q_b$ implies $P_b \cap Q_b = \emptyset$.
				By the definition of pairwise disjoint sets, we see that $A / \sim$ is pairwise disjoint.

				Let $c_0$ be an arbitrary element of $\bigcup_{P \in A / \sim} P$.
				By the definition of indexed union, we see that there exists a particular element $P_c \in A / \sim$ so that $c_0 \in P_c$.
				By the definition of set quotients, we see that there exists a particular element $c_1 \in A$ so that $c_0 \in \sim [c_1]$.
				Since $\sim$ is an equivalence relation on $A$, we see that $c_0 \in A$.
				By the definition of subsets, since $c_0 \in \bigcup_{P \in A / \sim} P$ implies $c_0 \in A$, we see that $\bigcup_{P \in A / \sim} P \subseteq A$.
				Let $d$ be an arbitrary element of $A$.
				By the definition of set quotients, we see that there exists a particular element $P_d \in A / \sim$ so that $P_d = \sim [d]$.
				By the definition of reflexive relations, we see that $d \in \sim [d]$.
				By substitution, we see that $d \in P_d$.
				By the definition of indexed union, since there exists a particular element $P_d \in A / \sim$ so that $d \in P_d$, we see that $d \in \bigcup_{P \in A / \sim} P$.
				By the definition of subsets, since $d \in A$ implies $d \in \bigcup_{P \in A / \sim} P$, we see that $A \subseteq \bigcup_{P \in A / \sim} P$.
				By the definition of set equality, we see that $\bigcup_{P \in A / \sim} P = A$.
				By the definition of exhaustive families of sets, we see that $A / \sim$ is exhaustive.

				By the definition of partitions, since $A / \sim$ contains no empty subsets, is pairwise disjoint, and is exhaustive, we see that $A / \sim$ is a partition.
				Therefore, we see that $F: E(A) \rightarrow \T_A$ is well-defined.
			\end{proof}

			\item
			Prove the function $G$ is well-defined by showing $G(\W) \in E(A)$.

			\begin{proof}
				We ask whether $G : \T_A \rightarrow E(A)$ is well-defined.
				Let $\W$ be an arbitrary partition in $\T_A$.
				Let $R_\W$ be the relation on $A$ induced by $W$.

				Let $a$ be an arbitrary element in $A$.
				By the definition of partitions, we see that there exists a particular element $P_a \in W$ so that $a \in P_a$.
				By the definition of $R_\W$, we see that $(a, a) \in R_\W$.
				By the definition of reflexive relations, we see that $R_\W$ is reflexive.

				Let $b_0$ and $b_1$ be arbitrary elements of $A$ so that $(b_0, b_1) \in R_\W$.
				By the definition of $R_\W$, we see that there exists a particular element $P_b \in W$ so that $b_0 \in P_b$ and $b_1 \in P_b$.
				By the definition of $R_\W$, we see that $(b_1, b_0) \in R_\W$.
				By the definition of symmetric relations, we see that $R_\W$ is symmetric.

				We ask whether $R_\W$ is transitive.
				Let $c_0, c_1$, and $c_2$ be arbitrary elements of $A$ so that $(c_0, c_1) \in R_\W$ and $(c_1, c_2) \in R_\W$.
				By the definition of $R_\W$, we see that there exists a unique element $P_c \in \W$ so that $c_0 \in P_c$ and so that $c_1 \in P_c$.
				By the definition of $R_\W$, we see that there exists a unique element $Q_c \in \W$ so that $c_1 \in Q_c$ and so that $c_2 \in Q_c$.
				By the definition of partitions, since $\W$ is pairwise disjoint, $c_1 \in P_c$, and $c_1 \in Q_c$, we see that $P_c = Q_c$.
				By substitution, we see that $c_2 \in P_c$.
				By the definition of $R_\W$, since $c_0 \in P_c$ and $c_2 \in P_c$, we see that $(c_0, c_2) \in R_\W$.
				By the definition of transitive relations, we see that $R_\W$ is transitive.

				By the definition of equivalence relations, since $R_\W$ is reflexive, symmetric, and transitive, we see that $R_\W$ is an equivalence relation.
				By substitution, since $R_\W = G(\W)$, we see that $G(\W)$ is an equivalence relation.
				By the definition of $E(A)$, we see that $G(\W) \in E(A)$.
				Therefore, we see that $G : \T_A \rightarrow E(A)$ is well-defined.
			\end{proof}

			\item
			Prove that $F \circ G = \id{\T_A}$.

			\begin{proof}
				We ask whether $F \circ G = \id{\T_A}$.
				By the definition of function composition, we see that $F \circ G : \T_A \rightarrow \T_A$.
				By the definition of the identity function, we see that $\id{\T_A} : \T_A \rightarrow \T_A$.
				We see that the domain and co-domain of $F \circ G$ and $\id{\T_A}$ are the same.
				Let $\W$ be an arbitrary element of $\T_A$.
				By the definition of the identity function, we see that $\id{\T_A}(\W) = \W$.
				By the definition of function composition and substitution, we see that
				$(F \circ G)(\W) = F(G(\W)) = F(R_\W) = A / R_\W$.

				Let $A_0$ be an arbitrary element of $A / R_\W$.
				By the definition of set quotients, we see that there exists a particular element $a_0 \in A$ so that $A_0 = R_\W [a_0]$.
				By the definition of partitions, since $\W$ is exhaustive, we see that there exists a particular element $P_0 \in \W$ so that $a_0 \in P_0$.
				By the definition of $R_\W$, we see that there exists a unique element $R_\W [a_0] \in \W$.
				By the definition of reflexive relations, we see that $a_0 \in R_\W [a_0]$.
				By the definition of partitions, since $\W$ is pairwise disjoint, $a_0 \in P_0$, $a_0 \in R_\W [a_0]$, we see that $P_0 = R_\W [a_0]$.
				By substitution, we see that $A_0 = P_0$.
				By substitution, we see that $A_0 \in \W$.
				By the definition of subsets, since $A_0 \in A / R_\W$ implies $A_0 \in \W$, we see that $A / R_\W \subseteq \W$.

				Let $A_1$ be an arbitrary element of $\W$.
				By the definition of partitions, since $\W$ contains no empty subsets, there exists a particular element $a_1 \in A_1$.
				By the definition of $\W$, since $A_1 \in \W$, we see that $a_1 \in A$.
				By the definition of $R_\W$, we see that there exists a unique element $R_\W [a_1] \in \W$.
				By the definition of reflexive relations, we see that $a_1 \in R_\W [a_1]$.
				By the definition of partitions, since $\W$ is pairwise disjoint, $a_1 \in A_1$, and $a_1 \in R_\W [a_1]$, we see that $A_1 = R_\W [a_1]$.
				By the definition of set quotients, we see that $A_1 \in A / R_\W$.
				By the definition of subsets, since $A_1 \in \W$ implies $A_1 \in A / R_\W$, we see that $\W \subseteq A / R_\W$.

				By the definition of set equality, we see that $A / R_\W = \W$.
				By substitution, we see that $(F \circ G)(\W) = \id{\T_A}(\W)$.
				Therefore, by the definition of function equality, since the domain and co-domain of $F \circ G$ and $\id{\T_A}$ are the same and  $(F \circ G)(\W) = \id{\T_A}(\W)$ for all elements $\W \in T_A$, we see that $(F \circ G) = \id{\T_A}$.
			\end{proof}

			\item
			Prove that $G \circ F = \id{E(A)}$.

			\begin{proof}
				We ask whether $G \circ F = \id{E(A)}$.
				By the definition of function composition, we see that $G \circ F : E(A) \rightarrow E(A)$.
				By the definition of the identity function, we see that $\id{E(A)} : E(A) \rightarrow E(A)$.
				We see that the domain and co-domain of $G \circ F$ and $\id{E(A)}$ are the same.
				Let $\sim$ be an arbitrary element of $E(A)$.
				By the definition of the identity function, we see that $\id{E(A)}(\sim) = \sim$.
				By the definition of function composition and substitution, we see that
				$(G \circ F)(\sim)  = G(F(\sim)) = G(A / \sim) = R_{A / \sim}$.

				Let $(a_0, a_1)$ be an arbitrary element of $\sim$.
				By the definition of equivalence classes, we see that $a_1 \in \sim [a_0]$.
				By the definition of reflexive relations, we see that $a_0 \in \sim [a_0]$.
				By the definition of set quotients, we see that $\sim [a_0] \in A / \sim$.
				By the definition of equivalence relations induced by partitions, we see that $(a_0, a_1) \in R_{A / \sim}$.
				By the definition of subsets, since $(a_0, a_1) \in \sim$ implies $(a_0, a_1) \in R_{A / \sim}$, we see that $\sim \subseteq R_{A / \sim}$.

				Let $(b_0, b_1)$ be an arbitrary element of $R_{A / \sim}$.
				By the definition of equivalence relations induced by partitions, there exists a particular element $B \in A / \sim$ so that $b_0 \in B$ and $b_1 \in B$.
				By the definition of set quotients, we see that there exists a particular element $b_2 \in A$ so that $B = \sim [b_2]$.
				By substitution, we see that $b_0 \in \sim [b_2]$ and $b_1 \in \sim [b_2]$.
				By the definition of equivalence classes, we see that $(b_2, b_0) \in \sim$ and $(b_2, b_1) \in \sim$.
				By the definition of symmetric relations, we see that $(b_0, b_2) \in \sim$.
				By the definition of transitive relations, we see that $(b_0, b_1) \in \sim$.
				By the definition of subsets, since $(b_0, b_1) \in R_{A / \sim}$ implies $(b_0, b_1) \in \sim$, we see that $R_{A / \sim} \subseteq \sim$.

				By the definition of set equality, we see that $R_{A / \sim} = \sim$.
				By substitution, we see that $(G \circ F)(\sim) = \id{E(A)}(\sim)$.
				Therefore, by the definition of function equality, since the domain and co-domain of $G \circ F$ and $\id{E(A)}$ are the same and  $(G \circ F)(\sim) = \id{E(A)}(\sim)$ for all elements $\sim \in E(A)$, we see that $(G \circ F) = \id{E(A)}$.
			\end{proof}
		\end{enumerate}
	\end{problem}

\end{document}